% =======================================================
% 4. BENCHMARKS
% =======================================================
\section{Execution-Validated Loop Benchmarks}
\label{sec:bench}

An evaluation of a parallelization oracle needs labels that do not come from any of the tools being compared. We therefore built two benchmarks whose labels are written under the legality definition of Section~\ref{sec:problem} and then checked by execution, and we keep the 200-file corpus of our preliminary study as an application-scale corpus without labels.

\subsection{HF-LoopBench}
HF-LoopBench is built from \nBenchBases{} hand-written \emph{base kernels}, each a small C function from one of \nBenchDomains{} application domains with exactly one \emph{target loop} and a documented ground-truth (GT) label. The kernels cover the idioms that decide parallelizability in practice: affine DOALL loops and stencils, scalar and array reductions, privatizable and \Code{linear} scalars, gathers through lookup tables or index arrays, disjoint but non-affine writes (bit-manipulated quantum-gate indices, compare-exchange networks, in-place reversals), permutation scatters, and on the serial side flow and anti dependences, scalar recurrences (filters, generators, hash chains), wavefronts, compaction counters, conflicting scatters, order-dependent library calls and early exits.

A generator instantiates every base kernel along three axes: the element type (up to two of \Code{double}, \Code{float}, \Code{int}, \Code{long long}); the data-access form, \textbf{G} (file-scope arrays), \textbf{P} (pointer parameters) or \textbf{R} (\Code{restrict}-qualified parameters); and, for single-input/single-output DOALL and stencil kernels, a \textbf{Q} variant in which the pointer version is \emph{called} with overlapping arguments (\Code{out = buf+1, in = buf}). The Q variant has the same function text as P but is serial, so it can only be classified correctly by looking at the call site. Each instance becomes a function in a domain file together with a test harness that initializes inputs deterministically, calls the kernel once and prints a 64-bit FNV hash of every output array, two order-sensitive sums and the return value. For every file the generator writes the clean source given to all tools (no directives, no labels), an \emph{OpenMP} version with the GT directive on every parallel target loop, and a \emph{race} version with a naive \Code{parallel for} on every canonical serial target loop. Table~\ref{tab:bench} summarizes the result. Instances of one base kernel are correlated, so all confidence intervals are computed with a cluster bootstrap over base kernels and $\tau$ is calibrated on a split by base kernel.

% generated by hf/report.py from results/metrics.json - do not edit by hand
\begin{table}[htbp]
\caption{Composition of HF-LoopBench (all numbers counted from the generated sources).}
\label{tab:bench}
\centering
\small
\resizebox{\linewidth}{!}{%
\begin{tabular}{lrrrrrr}
\toprule
Domain & Files & Lines & Base kernels & Target loops & GT par. & GT ser. \\
\midrule
bioinformatics & 2 & 1,018 & 10 & 33 & 21 & 12 \\
climate weather & 3 & 1,840 & 9 & 57 & 42 & 15 \\
cryptography & 2 & 895 & 9 & 29 & 15 & 14 \\
databases & 2 & 1,484 & 11 & 50 & 27 & 23 \\
finance & 3 & 2,256 & 12 & 75 & 42 & 33 \\
geometry graphics & 2 & 1,610 & 8 & 49 & 36 & 13 \\
graph analytics & 2 & 1,647 & 11 & 48 & 27 & 21 \\
image processing & 3 & 2,084 & 12 & 70 & 51 & 19 \\
linear algebra & 4 & 2,749 & 15 & 92 & 66 & 26 \\
machine learning & 4 & 2,787 & 14 & 89 & 66 & 23 \\
medical imaging & 2 & 1,247 & 6 & 38 & 24 & 14 \\
molecular dynamics & 2 & 1,692 & 9 & 51 & 42 & 9 \\
networking & 1 & 787 & 6 & 23 & 9 & 14 \\
quantum simulation & 2 & 1,176 & 6 & 36 & 30 & 6 \\
scientific misc & 3 & 1,727 & 10 & 59 & 42 & 17 \\
signal processing & 4 & 2,663 & 14 & 89 & 54 & 35 \\
sorting searching & 2 & 1,217 & 9 & 39 & 18 & 21 \\
stencil pde & 3 & 1,850 & 10 & 63 & 42 & 21 \\
\textbf{Total} & 46 & 30,729 & 181 & 990 & 654 & 336 \\
\bottomrule
\end{tabular}
}
\end{table}


\subsection{Dynamic validation of every label}
\label{sec:oracle}
All three versions of every file are compiled with Clang~21 at \Code{-O1 -ffp-contract=off}, so that floating-point evaluation is identical in the serial and OpenMP builds. The OpenMP version is run with 2, 4 and 8 threads under the schedules \Code{static}, \Code{static,1}, \Code{dynamic,1} and \Code{guided} (12 runs); a parallel label is accepted only if every run reproduces the serial output (bit-exact hash, or relative tolerance $10^{-8}$/$10^{-3}$ on the sums for \Code{double}/\Code{float} reductions). The race version is run with 4 and 8 threads under \Code{static,1} and \Code{dynamic,1}, three times each; a canonical serial label is \emph{confirmed} if at least one run deviates. Table~\ref{tab:oracle} shows the outcome: no parallel label was violated, and every canonical serial label of HF-LoopBench produced a wrong result under naive parallelization.

% generated by hf/report.py from results/metrics.json - do not edit by hand
\begin{table}[htbp]
\caption{Dynamic validation of the ground truth on the M1 Pro.}
\label{tab:oracle}
\centering
\small
\begin{tabular}{lr}
\toprule
Outcome & Loops \\
\midrule
GT-parallel loops verified (12 runs each) & 654 \\
GT-parallel loops violating their label & 0 \\
GT-serial loops whose naive parallelisation changed the output & 321 \\
GT-serial loops without an observed deviation & 0 \\
GT-serial loops not in OpenMP canonical form & 15 \\
\bottomrule
\end{tabular}
\end{table}


The harness also exposed a toolchain defect. At \Code{-O2}, the Homebrew build of Clang~21.1.1 produced wrong results for integer whole-array and array-section reductions even with a single thread (in a 15-line reproducer, \Code{hist[px[i]]++} with \Code{reduction(+:hist)} lost about 0.6\% of the increments under \Code{static} and \Code{guided} scheduling, while \Code{-O1}, \Code{-O0} and \Code{-O2 -fno-vectorize} were correct). The oracle therefore uses \Code{-O1}; in the \Code{-O2} performance build of Section~\ref{sec:rq-perf}, the \nPerfMismatch{} affected kernels---all integer array reductions---are excluded from the speedup statistics. The episode illustrates why a proposed directive, whether produced by a compiler, a model or a human, should be validated by differential execution before it is merged.

\subsection{HF-LoopBench-X: a loop-structure extension}
\label{sec:benchx}
HF-LoopBench is organized by application domain and dependence class. To test whether the conclusions carry over to loop \emph{shapes} that it covers thinly, we wrote \xBases{} further base kernels \emph{after} all rules, prompts and thresholds of the main study had been frozen. They include three \emph{level pairs} whose loop nest is textually identical but whose target loop differs (the spatial loop inside a time-step loop is parallel, the time-step loop itself is not; a column loop inside a vertical accumulation is parallel, the level loop is not; a row-diffusion versus a wavefront update of the same grid); triangular nests; stride-2 and step-4 loops; decrementing loops; an unsigned induction variable; walking-pointer loops (\Code{*q++ = *p++ * s}, legal with \Code{linear(p,q:1)}); \Code{collapse(2)}; if/else, masked and \Code{continue} bodies; product and nested reductions; a privatizable temporary array; disjoint in-place updates (a stride-7 modulo interleaver over a power-of-two buffer, a half-buffer shift, pair swaps, a permuted update); and subtle serial loops (a cross-statement flow dependence, a distance-2 recurrence, an anti dependence with a run-time offset, a conditional recurrence, an 8-slot ring buffer, an arg-min, a linked-list walk). Besides the overlapping-call variant Q, every eligible kernel also gets a variant \textbf{H} that is called with two \emph{disjoint} halves of one buffer and therefore remains parallel; it measures how often the conservative call-site guard vetoes a legal loop. The generator, harness and oracle are those of HF-LoopBench. Of the \xLoops{} target loops (\xPar{} parallel, \xSer{} serial), all \xVerified{} parallel labels were verified and \xManifested{} serial labels produced a deviation; the remaining \xNotManifested{} are the arg-min loop, which is serial under our definition because OpenMP has no built-in reduction that returns both the minimum and its index, but whose race did not change the result in 12 runs.

\subsection{Application-scale corpus}
The 200 C++ files (\fieldLines{} lines) of our preliminary study---trading, computer vision, cryptography, finance, healthcare, machine learning, networking, quantum computing, scientific computing and weather/climate code using classes, STL containers and helper functions---have no ground truth. We use them to measure how HF-Static scales and how many loops of realistic code would be routed to the model (Section~\ref{sec:rq-field}).
